% year: תשפ"ד
% semester: ב
% moed: ב
% number: 2
% categories: continuity, derivatives
% source: exams/מבחנים/תשפד/מועד ב תשפד סמסטר ב.pdf

\begin{question}
לכל $a,b\in\R$ נגדיר את הפונקציה הבאה:
\[ f(x)=\begin{cases} \sin(ax) & x<0\\ b-2 & x=0\\ x^2\ln(x) & 0>x\end{cases} \]
\begin{enumerate}
  \item עבור אילו $a,b\in\R$ הפונקציה $f(x)$ רציפה בנקודה $x=0$?
  \item האם קיימים $a,b\in\R$ עבורם הפונקציה $f(x)$ גזירה בנקודה $x=0$?
\end{enumerate}
\end{question}

\begin{hint}
חשבו את הגבולות החד-צדדיים ב-$0$. לגבול $\lim_{x\to0^+}x^2\ln x$ (וגם $x\ln x$) השתמשו בכלל לופיטל אחרי כתיבה כמנה.
\end{hint}

\begin{hint}
גזירות גוררת רציפות. חשבו את הנגזרות החד-צדדיות לפי ההגדרה.
\end{hint}

\begin{solution}
\textbf{הערה:} בתנאי של השורה השלישית נכתב "$0>x$", וזו כמובן טעות דפוס — השורה הראשונה כבר מכסה את $x<0$, וב-$x<0$ הביטוי $\ln x$ אינו מוגדר. הכוונה היא $0<x$, ונפתור בהתאם.
\begin{enumerate}
  \item \textbf{גבול משמאל:} $\lim_{x\to0^-}\sin(ax)=\sin0=0$ (רציפות הסינוס), לכל $a$.

  \textbf{גבול מימין:} $\lim_{x\to0^+}x^2\ln x$ הוא מהצורה $0\cdot(-\infty)$. נכתוב כמנה ונשתמש בכלל לופיטל:
  \[ \lim_{x\to0^+}\frac{\ln x}{x^{-2}}=\lim_{x\to0^+}\frac{1/x}{-2x^{-3}}=\lim_{x\to0^+}\left(-\frac{x^2}{2}\right)=0. \]
  לכן הגבול של $f$ ב-$0$ קיים ושווה $0$ לכל $a$, ו-$f$ רציפה ב-$0$ אם ורק אם $f(0)=b-2=0$.
  \textbf{תשובה:} $f$ רציפה ב-$0$ עבור $b=2$ ו-$a\in\R$ כלשהו.
  \item אם $f$ גזירה ב-$0$ היא בפרט רציפה שם, ולכן חייב להתקיים $b=2$, כלומר $f(0)=0$. נחשב את הנגזרות החד-צדדיות לפי ההגדרה:
  \[ f'_-(0)=\lim_{h\to0^-}\frac{\sin(ah)-0}{h}=a\lim_{h\to0^-}\frac{\sin(ah)}{ah}=a \]
  (עבור $a=0$ המנה זהותית $0$, וגם אז הגבול $0=a$), ו-
  \[ f'_+(0)=\lim_{h\to0^+}\frac{h^2\ln h-0}{h}=\lim_{h\to0^+}h\ln h=\lim_{h\to0^+}\frac{\ln h}{1/h}\overset{\text{לופיטל}}{=}\lim_{h\to0^+}\frac{1/h}{-1/h^2}=\lim_{h\to0^+}(-h)=0. \]
  $f$ גזירה ב-$0$ אם ורק אם שתי הנגזרות החד-צדדיות שוות, כלומר $a=0$.
  \textbf{תשובה:} כן — $f$ גזירה ב-$0$ בדיוק עבור $a=0,\ b=2$, ואז $f'(0)=0$.
\end{enumerate}
\end{solution}
